% Part: normal-modal-logic
% Chapter: completeness
% Section: frame-completeness

\documentclass[../../../include/open-logic-section]{subfiles}

\begin{document}

\olfileid{nml}{com}{fra}

\olsection{د چوکاټ بشپړتيا}

د \Log{K} د بشپړتيا قضيه نورو مودالي نظامونو ته غځېدے شي، کله چې وښيو
چې د يوه ټاکلي منطق کانوني مدل د چوکاټ اړوند خاصيت لري۔

\begin{thm}\ollabel{thm:completeframeprops}
  که يو عادي مودالي منطق $\Sigma$ د \olref{tab:correspondencetable} د کيڼ لوري
  له !!{formula}s څخه يو ولري،
  نو د~$\Sigma$ کانوني مدل د ښي لوري اړوند خاصيت لري۔
\end{thm}

\begin{table}[htp]
  \centering
    \begin{tabular}{| l || l |}
      \hline
      \emph{که $\Sigma$ دا ولري \dots} & \emph{\dots نو د $\Sigma$ کانوني
        مدل دا خاصيت لري:} \\
      \hline \hline
      \Ax{D}:\;\; $\Box!A \lif \Diamond !A$ & \quad مسلسل؛ \\
      \hline
      \Ax{T}:\;\; $\Box!A \lif !A$ &\quad  انعکاسي؛\\
      \hline
      \Ax{B}: \;\;$!A \lif \Box\Diamond!A$ &\quad  متناظر؛ \\
      \hline
      \Ax{4}: \;\; $\Box!A \lif \Box\Box!A$ & \quad متعدي؛ \\
      \hline 
      \Ax{5}: \;\; $\Diamond !A \lif \Box\Diamond!A$& \quad اقليدسي۔\\
      \hline
    \end{tabular}
    \caption{د مطابقت بنسټيز واقعيتونه۔}
    \ollabel{tab:correspondencetable}
\end{table}

\begin{proof}
  دا هر يو په خپل وار اخلو۔

  فرض کړو چې $\Sigma$، \Ax{D} لري، او $\Delta \in W^\Sigma$ واخلو؛ بايد
  وښيو چې داسې يو $\Delta'$ شته چې $R^\Sigma
  \Delta\Delta'$ وي۔ بس ده چې وښيو $\Box^{-1}\Delta$
  $\Sigma$-سازګار دے، ځکه بيا د لينډنباوم لمې له مخې يو
  بشپړ $\Sigma$-سازګار سټ $\Delta' \supseteq
  \Box^{-1}\Delta$ شته؛ او د $R^\Sigma$ د تعريف
  \iftag{prvBox}{}{او \olref[mod]{lem:box-iff-diamond} }له مخې
  $R^\Sigma \Delta\Delta'$ لرو۔ نو د تضاد لپاره فرض کړو چې
  $\Box^{-1}\Delta$ $\Sigma$-سازګار \emph{نۀ} دے، يعنې
  $\Box^{-1}\Delta \Proves[\Sigma] \lfalse$۔ د \olref[mod]{lem:box2} له مخې،
  $\Delta \Proves[\Sigma] \Box\lfalse$؛ او چونکې $\Sigma$،
  \Ax{D} لري، نو $\Delta \Proves[\Sigma] \Diamond\lfalse$ هم۔ خو $\Sigma$
  عادي دے، نو $\Sigma \Proves \lnot\Diamond \lfalse$
  (\olref[prf][nor]{prop:notDiamondBot})؛ له دې امله $\Delta
  \Proves[\Sigma] \lnot\Diamond \lfalse$ هم، چې د~$\Delta$
  د سازګارۍ خلاف دے۔

  اوس فرض کړو چې $\Sigma$، \Ax{T} لري، او $\Delta \in
  W^\Sigma$ واخلو۔ غواړو وښيو چې $R^\Sigma \Delta\Delta$،
  يعنې\iftag{prvBox}{}{ د \olref[mod]{lem:box-iff-diamond} له مخې،}
  $\Box^{-1}\Delta \subseteq \Delta$۔ خو که $\Box!A \in \Delta$ وي، نو
  د \Ax{T} له مخې $!A \in \Delta$ هم، لکه چې غوښتل شوي وو۔

  اوس فرض کړو چې $\Sigma$، \Ax{B} لري، او د $\Delta$، $\Delta' \in W^\Sigma$
  لپاره $R^\Sigma \Delta\Delta'$ فرض کړو۔ بايد وښيو
  چې $R^\Sigma \Delta'\Delta$، يعنې\iftag{prvBox}{
    $\Box^{-1}\Delta' \subseteq \Delta$۔ د
    \olref[mod]{lem:box-iff-diamond} له مخې، دا له دې سره همارز دے چې}{}
  $\Diamond\Delta \subseteq \Delta'$۔ نو $!A \in \Delta$ فرض کړو۔ د
  \Ax{B} له مخې، $\Box\Diamond!A \in \Delta$ هم۔ له هغې فرضيې چې
  $R^\Sigma \Delta\Delta'$\iftag{prvBox}{}{ او
    \olref[mod]{lem:box-iff-diamond}}، لرو چې $\Box^{-1}\Delta
  \subseteq \Delta'$؛ نو $\Diamond!A \in \Delta'$، لکه چې غوښتل شوي وو۔

  اوس فرض کړو چې $\Sigma$، \Ax{4} لري، او $R^\Sigma
  \Delta_1\Delta_2$ او $R^\Sigma \Delta_2\Delta_3$ فرض کړو۔ بايد وښيو
  چې $R^\Sigma \Delta_1\Delta_3$۔ له فرضيې\iftag{prvBox}{}{ او
    \olref[mod]{lem:box-iff-diamond}} دواړه $\Box^{-1}\Delta_1
  \subseteq \Delta_2$ او $\Box^{-1}\Delta_2 \subseteq \Delta_3$ لرو۔ د
  $R^\Sigma \Delta_1\Delta_3$ د ښودلو لپاره بس ده چې
  $\Box^{-1}\Delta_1 \subseteq \Delta_3$ وښيو۔ نو $!B \in
  \Box^{-1}\Delta_1$ واخلو، يعنې $\Box!B \in \Delta_1$۔ د \Ax{4} له مخې،
  $\Box\Box!B \in \Delta_1$ هم؛ او له فرضيې لومړے دا ترلاسه کوو چې
  $\Box!B \in \Delta_2$، او بيا دا چې $!B \in \Delta_3$، لکه چې
  غوښتل شوي وو۔

  اوس فرض کړو چې $\Sigma$، \Ax{5} لري، او $R^\Sigma
  \Delta_1\Delta_2$ او $R^\Sigma \Delta_1\Delta_3$ فرض کړو۔ بايد وښيو
  چې $R^\Sigma \Delta_2\Delta_3$۔ لومړۍ فرضيه
  $\Box^{-1}\Delta_1 \subseteq \Delta_2$\iftag{prvBox}{}{ د
    \olref[mod]{lem:box-iff-diamond} له مخې} راکوي، او دويمه فرضيه
  له دې سره همارزه ده چې $\Diamond\Delta_3 \subseteq \Delta_1$\iftag{prvBox}{،
    د \olref[mod]{lem:box-iff-diamond} له مخې}{}۔ د $R^\Sigma
  \Delta_2\Delta_3$ د ښودلو لپاره\iftag{prvBox}{، د
    \olref[mod]{lem:box-iff-diamond} له مخې، بس ده چې}{ بايد}
  $\Diamond\Delta_3 \subseteq \Delta_2$ وښيو۔ نو $\Diamond!A \in
  \Diamond\Delta_3$ واخلو، يعنې $!A \in \Delta_3$۔ د دويمې فرضيې له مخې
  $\Diamond!A \in \Delta_1$؛ او د \Ax{5} له مخې، $\Box\Diamond!A \in
  \Delta_1$ هم۔ خو اوس لومړۍ فرضيه $\Diamond!A \in
  \Delta_2$ راکوي، لکه چې غوښتل شوي وو۔
  \end{proof}

د يوې نتيجې په توګه د څو نظامونو د بشپړتيا پايلې ترلاسه کوو۔
د بېلګې په توګه، پوهېږو چې $\Log{S5} = \Log{KT5} =
\Log{KTB4}$ د ټولو انعکاسي اقليدسي مدلونو د ټولګي په نسبت بشپړ دے؛
دا عين د ټولو انعکاسي، متناظرو او متعدي مدلونو ټولګے دے۔

\begin{thm}\ollabel{thm:generaldet}
  پرېږدئ چې $\mClass{C}_\Ax{D}$، $\mClass{C}_\Ax{T}$،
  $\mClass{C}_\Ax{B}$، $\mClass{C}_\Ax{4}$ او
  $\mClass{C}_\Ax{5}$ په ترتيب د ټولو مسلسل، انعکاسي،
  متناظرو، متعدي او اقليدسي مدلونو ټولګي وي۔ بيا د
  \Ax{D}، \Ax{T}، \Ax{B}، \Ax{4} او \Ax{5} له منځه د هر ډول
  سکيماګانو $!A_1$، \dots، $!A_n$ لپاره، نظام
  $\Log{K}!A_1 \dots !A_n$ د مدلونو د هغه ټولګي له خوا ټاکل کېږي چې
  $\mClass{C} = \mClass{C}_{!A_1} \cap \dots
  \cap \mClass{C}_{!A_n}$ دے۔
\end{thm}

\begin{prop}
  پرېږدئ چې $\Sigma$ يو عادي مودالي منطق وي؛ بيا:
  \begin{enumerate}
  \item\ollabel{prop:anotherfive-a}%
    که $\Sigma$ د $\Diamond!A \lif \Box
    !A$ سکيما ولري، نو د $\Sigma$ کانوني مدل جزوي تابعي دے۔
  \item که $\Sigma$ د $\Diamond!A \liff \Box
    !A$ سکيما ولري، نو د $\Sigma$ کانوني مدل تابعي دے۔
  \item که $\Sigma$ د $\Box\Box!A \lif \Box
    !A$ سکيما ولري، نو د $\Sigma$ کانوني مدل کمزورې ګڼه لري۔
  \end{enumerate}
(د چوکاټ د دغو خاصيتونو د تعريفونو لپاره \olref[frd][acc]{tab:anotherfive}
  وګورئ)۔
\end{prop}

\begin{proof}
  \begin{enumerate}
  \item فرض کړو چې $\Sigma$ د $\Diamond !A \lif
    \Box !A$ سکيما لري۔ د دې د ښودلو لپاره چې $R^\Sigma$ جزوي تابعي دے، بايد
    ثابت کړو چې د هر $\Delta_1$، $\Delta_2$، $\Delta_3 \in
    W^\Sigma$ لپاره، که $R^\Sigma \Delta_1\Delta_2$ او $R^\Sigma
    \Delta_1\Delta_3$ وي، نو $\Delta_2=\Delta_3$۔ چونکې $R^\Sigma
    \Delta_1\Delta_2$، نو $\Box^{-1}\Delta_1 \subseteq \Delta_2$ لرو؛
    او چونکې $R^\Sigma \Delta_1\Delta_3$، نو $\Box^{-1}\Delta_1
    \subseteq \Delta_3$ هم\iftag{prvBox}{}{، دواړه د
      \olref[mod]{lem:box-iff-diamond} له مخې}۔ د $\Delta_2=\Delta_3$ عين‌والے
    هله ترلاسه کېږي چې دواړه شمولونه، $\Delta_2 \subseteq \Delta_3$ او $\Delta_3 \subseteq
    \Delta_2$، ثابت کړو۔ د لومړي شمول لپاره، $!A \in \Delta_2$ واخلو؛ بيا
    $\Diamond!A \in \Delta_1$، او د سکيما او د $\Delta_1$ د اشتقاقي تړلتيا
    له مخې $\Box!A \in \Delta_1$ هم؛ بيا له هغې فرضيې چې
    $R^\Sigma \Delta_1\Delta_3$، $!A \in \Delta_3$۔ دويم
    شمول هم په همدغه ډول دے۔

  \item دا سمدستي له~\olref{prop:anotherfive-a} برخې
    او په \olref{thm:completeframeprops} کښې د مسلسل‌والي له ثبوته راځي۔

  \item فرض کړو چې $\Sigma$ د $\Box\Box!A \lif \Box!A$ سکيما لري۔
    د دې د ښودلو لپاره چې $R^\Sigma$ کمزورې ګڼه لري، $R^\Sigma
    \Delta_1\Delta_2$ واخلو۔ بايد وښيو چې داسې يو بشپړ
    $\Sigma$-سازګار سټ $\Delta_3$ شته چې $R^\Sigma
    \Delta_1\Delta_3$ او $R^\Sigma \Delta_3\Delta_2$ وي۔ پرېږدئ چې:
    \[
    \Gamma = \Box^{-1}\Delta_1 \cup \Diamond\Delta_2.
    \]
    بس ده چې وښيو $\Gamma$ $\Sigma$-سازګار دے، ځکه بيا
    د لينډنباوم لمې له مخې يو بشپړ
    $\Sigma$-سازګار سټ~$\Delta_3$ ته غځېدے شي، داسې چې $\Box^{-1}\Delta_1
    \subseteq \Delta_3$ او $\Diamond\Delta_2 \subseteq \Delta_3$،
    يعنې $R^\Sigma \Delta_1\Delta_3$\iftag{prvBox}{}{ (د
      \olref[mod]{lem:box-iff-diamond} له مخې)} او $R^\Sigma
    \Delta_3\Delta_2$\iftag{prvBox}{ (د
      \olref[mod]{lem:box-iff-diamond} له مخې)}{}۔

    د تضاد لپاره فرض کړو چې $\Gamma$ سازګار نۀ دے۔ بيا
    داسې !!{formula}s $\Box!A_1$، \dots، $\Box!A_n \in \Delta_1$
    او $!B_1$، \dots،~$!B_m \in \Delta_2$ شته چې \[!A_1, \dots,
    !A_n, \Diamond!B_1, \dots, \Diamond!B_m \Proves[\Sigma]
    \lfalse.\] چونکې $\Diamond (!B_1 \land \dots \land !B_m) \to
    (\Diamond!B_1 \land \dots \land \Diamond!B_m)$ په
    هر عادي مودالي منطق کښې !!{derivable} دے، نو لاندې استدلال کوو، چې د~$\Delta_2$
    له سازګارۍ سره تضاد لري:
    \begin{align*}
      !A_1, \dots, !A_n, & \Diamond!B_1, \dots,
      \Diamond!B_m \Proves[\Sigma] \lfalse \\
      !A_1, \dots,!A_n & \Proves[\Sigma] (\Diamond!B_1 \land
      \dots \land \Diamond!B_m) \lif \lfalse\\
      & \qquad\text{د استنتاج د قضيې له مخې}\\
      & \qquad\text{\olref[prf][prp]{prop:derivabilityfacts}\olref[prf][prp]{prop:derivabilityfacts-deduction}، او \Taut} \\
      !A_1, \dots,!A_n & \Proves[\Sigma]
      \Diamond(!B_1 \land
      \dots \land !B_m) \lif \lfalse\\
      & \qquad \text{چونکې $\Sigma$ عادي دے} \\
      !A_1, \dots,!A_n & \Proves[\Sigma]
      \lnot\Diamond (!B_1 \land \dots \land !B_m)\\
      & \qquad\text{د \PL له مخې} \\
      !A_1, \dots,!A_n & \Proves[\Sigma]
      \Box\lnot (!B_1 \land \dots \land !B_m)\\
      & \qquad\text{$\Box\lnot$ د $\lnot\Diamond$ په معنا} \\
      \Box!A_1, \dots,\Box!A_n
      & \Proves[\Sigma] \Box\Box \lnot (!B_1 \land \dots \land !B_m)\\
      & \qquad\text{د \olref[mod]{lem:box1} له مخې} \\
      \Box!A_1, \dots,\Box!A_n
      & \Proves[\Sigma]
      \Box\lnot (!B_1 \land \dots \land !B_m)\\
      &\qquad\text{د $\Box\Box!A \lif \Box!A$ سکيما له مخې} \\
      \Delta_1 & \Proves[\Sigma]
      \Box\lnot (!B_1 \land \dots \land !B_m)\\
      &\qquad\text{د يو‌رنګه‌والي له مخې، \olref[prf][prp]{prop:derivabilityfacts}\olref[prf][prp]{prop:derivabilityfacts-monotonicity}} \\
      & \Box\lnot (!B_1 \land \dots \land !B_m) \in \Delta_1\\
      &\qquad\text{د اشتقاقي تړلتيا له مخې}; \\
      & \lnot (!B_1 \land \dots \land !B_m) \in \Delta_2\\
      &\qquad \text{چونکې }
      R^\Sigma \Delta_1\Delta_2.
    \end{align*}
  \end{enumerate}
\end{proof}

د دغو بېلګو پر بنسټ ښايي داسې وګڼل شي چې د مودالي منطق هر
نظام~$\Sigma$ \emph{بشپړ} دے، په دې معنا چې هر هغه فارمول ثابتوي چې په هر
هغه چوکاټ کښې اعتبار لري چې د $\Sigma$ هره قضيه پکښې اعتبار لري۔
له بده مرغه، ډېر نظامونه شته چې په دې معنا بشپړ نۀ دي۔

\end{document}
